Este anexo detalla las instrucciones necesarias para clonar el repositorio, configurar el entorno de desarrollo local y ejecutar las pruebas de \textbf{Sky Apartments}.

\section{Configuración del Repositorio}
Para clonar el repositorio del proyecto en el equipo local, se debe ejecutar el siguiente comando:

\begin{lstlisting}[language=bash]
git clone https://github.com/codeurjc-students/2025-sky-apartments.git
cd 2025-sky-apartments
\end{lstlisting}

\section{Ejecución con Docker Compose (Recomendado)}
La forma más sencilla de ejecutar la aplicación completa con todos sus servicios es utilizando Docker Compose.

\subsection{Versión Estable}
\begin{lstlisting}[language=bash]
cd docker
docker compose up
\end{lstlisting}
Este comando realizará las siguientes acciones:
\begin{itemize}
    \item Descargar la última imagen estable desde DockerHub.
    \item Iniciar contenedores MySQL independientes para cada microservicio.
    \item Iniciar MinIO para el almacenamiento de objetos.
    \item Iniciar Eureka Server para el descubrimiento de servicios.
    \item Iniciar todos los microservicios de la lógica de negocio (Usuarios, Apartamentos, Reservas, Reseñas).
    \item Iniciar el API Gateway con configuración HTTPS.
    \item Iniciar Jaeger para la trazabilidad distribuida.
    \item Desplegar el \textit{frontend} accesible en \url{https://localhost/}.
\end{itemize}

\subsection{Versión de Desarrollo}
Para probar los últimos cambios integrados en la rama principal (\texttt{main}), se puede utilizar la versión de desarrollo:
\begin{lstlisting}[language=bash]
cd docker
docker compose -f docker-compose-dev.yml up
\end{lstlisting}

\subsection{Accesos a los servicios}
Una vez levantada la infraestructura, los distintos servicios son accesibles a través de las siguientes rutas:
\begin{itemize}
    \item \textbf{Frontend:} \url{https://localhost/}
    \item \textbf{API Gateway:} \url{https://localhost/api/}
    \item \textbf{Panel de Eureka:} \url{http://localhost:8761/}
    \item \textbf{Interfaz de Jaeger:} \url{http://localhost:16686/}
    \item \textbf{Consola de MinIO:} \url{http://localhost:9001/}
\end{itemize}

\section{Ejecución Manual (Entorno de Desarrollo Local)}
Para modificar el código y realizar desarrollo local sin empaquetar en contenedores Docker la propia aplicación, se deben seguir estos pasos.

\subsection{Requisitos Previos}
\begin{itemize}
    \item \textbf{Java:} Versión 17.
    \item \textbf{Node.js:} Versión 20.18.0.
    \item \textbf{MySQL:} Versión 8.0 (mediante Docker o instalación nativa).
    \item \textbf{MinIO:} Para almacenamiento de objetos (mediante Docker).
    \item \textbf{Maven Wrapper:} Incluido en el propio repositorio.
\end{itemize}

\subsection{Configuración de Bases de Datos}
Se deben arrancar los contenedores de MySQL para cada microservicio ejecutando:

\begin{lstlisting}[language=bash]
# Base de datos del Servicio de Usuarios (Puerto 3307)
docker run --name mysql-user -e MYSQL_ROOT_PASSWORD=root -e MYSQL_DATABASE=user_db -e MYSQL_USER=user -e MYSQL_PASSWORD=password -p 3307:3306 -d mysql:8.0

# Base de datos del Servicio de Apartamentos (Puerto 3308)
docker run --name mysql-apartment -e MYSQL_ROOT_PASSWORD=root -e MYSQL_DATABASE=apartment_db -e MYSQL_USER=user -e MYSQL_PASSWORD=password -p 3308:3306 -d mysql:8.0

# Base de datos del Servicio de Reservas (Puerto 3309)
docker run --name mysql-booking -e MYSQL_ROOT_PASSWORD=root -e MYSQL_DATABASE=booking_db -e MYSQL_USER=user -e MYSQL_PASSWORD=password -p 3309:3306 -d mysql:8.0

# Base de datos del Servicio de Resenas (Puerto 3310)
docker run --name mysql-review -e MYSQL_ROOT_PASSWORD=root -e MYSQL_DATABASE=review_db -e MYSQL_USER=user -e MYSQL_PASSWORD=password -p 3310:3306 -d mysql:8.0
\end{lstlisting}

\subsection{Configuración de MinIO}
\begin{lstlisting}[language=bash]
docker run --name minio \
  -e MINIO_ROOT_USER=minioadmin \
  -e MINIO_ROOT_PASSWORD=minioadmin \
  -p 9000:9000 \
  -p 9001:9001 \
  -d minio/minio server /data --console-address ":9001"
\end{lstlisting}

\subsection{Ejecución del Backend (\textit{Microservicios})}
Es necesario abrir un terminal independiente para cada servicio e iniciarlos en el siguiente orden estricto:

\begin{lstlisting}[language=bash]
cd backend

# 1. Iniciar Eureka Server
cd eureka-server && mvn spring-boot:run

# 2. Iniciar Servicio de Usuarios
cd user && mvn spring-boot:run

# 3. Iniciar Servicio de Apartamentos
cd apartment && mvn spring-boot:run

# 4. Iniciar Servicio de Reservas
cd booking && mvn spring-boot:run

# 5. Iniciar Servicio de Resenas
cd review && mvn spring-boot:run

# 6. Iniciar API Gateway
cd api-gateway && mvn spring-boot:run
\end{lstlisting}

\subsection{Ejecución del Frontend (\textit{Angular})}
En un nuevo terminal:
\begin{lstlisting}[language=bash]
cd frontend
npm install
npm start # O alternativamente: ng serve
\end{lstlisting}
El \textit{frontend} estará disponible en \url{http://localhost:4200}. Las peticiones a \texttt{/api} se redirigirán automáticamente al API Gateway mediante la configuración del \textit{proxy} de desarrollo.

\section{Herramientas de Desarrollo Adicionales}
\begin{itemize}
    \item \textbf{IDE:} Visual Studio Code (tanto para los microservicios en Java como para Angular).
    \item \textbf{Postman:} Para interactuar con la API REST a través del Gateway. El repositorio incluye una colección con ejemplos de peticiones y datos de prueba en la ruta: \texttt{backend/postman/Sky\_Apartments.postman\_collection.json}.
\end{itemize}

\section{Ejecución de Pruebas Automáticas}

\subsection{Pruebas del Backend}

\subsubsection*{Pruebas Unitarias}
\begin{lstlisting}[language=bash]
cd backend

# Ejecutar pruebas en todos los microservicios
mvn test -Dtest='*UnitTest'

# Ejecutar pruebas en un microservicio especifico (Ej: apartment)
mvn test -Dtest='*UnitTest' -pl apartment
\end{lstlisting}

\subsubsection*{Pruebas de Integración}
\begin{lstlisting}[language=bash]
cd backend

# Ejecutar pruebas en todos los microservicios
mvn test -Dtest='*IntegrationTest'

# Ejecutar pruebas en un microservicio especifico
mvn test -Dtest='*IntegrationTest' -pl apartment
\end{lstlisting}

\subsubsection*{Pruebas End-to-End (E2E)}
Las pruebas E2E requieren que la infraestructura esté en funcionamiento:
\begin{lstlisting}[language=bash]
# Iniciar infraestructura de pruebas
cd docker
docker compose -f docker-compose-testing.yml up -d

# Esperar a que los servicios esten listos en Eureka y ejecutar:
cd ../backend
mvn test -Dtest='*e2eTest' -pl apartment
mvn test -Dtest='*e2eTest' -pl user

# Detener infraestructura al finalizar
cd ../docker
docker compose -f docker-compose-testing.yml down -v
\end{lstlisting}

\subsection{Pruebas del Frontend}
\begin{lstlisting}[language=bash]
cd frontend

# Pruebas Unitarias
npm run test:unit

# Pruebas de Integracion (Requiere backend en ejecucion)
npm run test:integration

# Pruebas E2E con Playwright (Requiere stack completo)
npm run test:e2e

# Ver el informe de resultados de Playwright
npx playwright show-report
\end{lstlisting}

\section{Construcción de Imágenes Docker}
Para compilar y construir la imagen de Docker localmente:
\begin{lstlisting}[language=bash]
cd docker
docker build -t skyapartments:local .
docker compose -f docker-compose-local.yml up
\end{lstlisting}

\section{Publicación y Gestión de Versiones (\textit{Releases})}
El proyecto implementa un flujo de \textbf{entrega continua} automatizado mediante GitHub Actions. La publicación de una nueva versión se activa manualmente al crear una \textit{Release} en GitHub, momento en el que se dispara automáticamente la \textit{pipeline} de CD. El paso final a producción (AWS EC2) requiere intervención manual del administrador. El proceso completo es el siguiente:

\begin{enumerate}
    \item \textbf{Actualización manual de versiones:}
    \begin{itemize}
        \item Archivo \texttt{backend/pom.xml} (Ej. \texttt{0.1.0}).
        \item Archivo \texttt{frontend/package.json} (Ej. \texttt{0.1.0}).
        \item Archivo \texttt{docker/docker-compose.yml} (Ej. \texttt{0.1}).
    \end{itemize}
    \item \textbf{Creación de la Release en GitHub:} Se genera una nueva \textit{Release} 
    asociándole una etiqueta (\textit{tag}) coincidente con la versión (Ej. \texttt{0.1}). 
    Este evento dispara automáticamente la \textit{pipeline} de entrega continua.
    \item \textbf{Ejecución automática de la \textit{pipeline}:} GitHub Actions construye 
    la imagen, actualiza la etiqueta \texttt{latest}, publica la imagen en DockerHub y sube 
    el archivo \texttt{docker-compose.yml} como artefacto OCI.
    \item \textbf{Preparación del próximo ciclo:} Se actualizan los archivos \texttt{pom.xml} 
    y \texttt{package.json} a la siguiente versión \texttt{SNAPSHOT} 
    (Ej. \texttt{0.2.0-SNAPSHOT}).
\end{enumerate}

\section{Acceso a Imágenes Publicadas}
Las imágenes generadas por la \textit{pipeline} están disponibles públicamente en DockerHub bajo las siguientes etiquetas:
\begin{itemize}
    \item \textbf{Versión específica:} \texttt{eloydsdlh/apartments-app:1.0}
    \item \textbf{Última versión estable:} \texttt{eloydsdlh/apartments-app:latest}
    \item \textbf{Última versión de desarrollo:} \texttt{eloydsdlh/apartments-app:dev}
\end{itemize}